\clearpage
\sectionkicker{Source-backed technical notes}
\subsection*{Frontier Models — 7月の世代交代: Technical Notes}
この欄は記事本文で圧縮した一次資料上の情報を、比較・再検証しやすい形へ展開したものである。新しい外部情報は追加せず、Selection済みEvidenceのchronology、normalized claim、limitations、source URLのみを再配置する。
\medskip
\subsection*{Theme at a glance}
\addcontentsline{toc}{subsection}{Theme at a glance}
\begin{center}
\footnotesize
\begin{tabularx}{\linewidth}{@{}p{0.20\linewidth}p{0.10\linewidth}p{0.14\linewidth}X@{}}
\toprule
Artifact & Role & Type & Objective chronology \\
\midrule
GPT-5.6 family & PRIMARY & MODEL & 2026-07-09 (MODEL\_RELEASE); 2026-07-30 (MODEL\_UPDATE) \\
Claude Opus 5 & PRIMARY & MODEL & 2026-07-24 (MODEL\_RELEASE) \\
DeepSeek V4-Flash July API lifecycle & PRIMARY & MODEL\_UPDATE & 2026-07-24 (API\_DEPRECATION); 2026-07-31 (MODEL\_UPDATE) \\
Gemini API July model lifecycle & PRIMARY & MODEL\_UPDATE & 2026-07-21 (MODEL\_GA); 2026-07-30 (API\_PREVIEW) \\
Kimi K3 & PRIMARY & OPEN\_WEIGHT & 2026-07-16 (MODEL\_RELEASE) \\
\bottomrule
\end{tabularx}
\end{center}
\normalsize

\begin{technicalnote}{GPT-5.6 family}{PRIMARY}
\begin{tabularx}{\linewidth}{@{}>{\bfseries}p{0.22\linewidth}X@{}}
Organization & OpenAI \\
Artifact type & MODEL \\
Chronology & 2026-07-09 (MODEL\_RELEASE); 2026-07-30 (MODEL\_UPDATE) \\
\end{tabularx}
\smallskip
{\bfseries 一次資料から整理したtechnical points}
\begin{itemize}[leftmargin=1.5em,itemsep=0.35em]
\item \textbf{一次情報で確認できる事実}: OpenAI launched the GPT-5.6 family for general availability on July 9, 2026 with Sol, Terra, and Luna tiers across ChatGPT, Codex, and the API.
\item \textbf{一次情報で確認できる事実}: The release introduced Programmatic Tool Calling in the Responses API and a multi-agent beta; OpenAI's `ultra` setting coordinates parallel agents for demanding tasks.
\item \textbf{一次情報で確認できる事実}: OpenAI later reduced GPT-5.6 Luna pricing by 80\% and Terra pricing by 20\% on July 30.
\item \textbf{Vendor claim}: Performance-per-dollar and benchmark superiority claims in the launch and follow-up articles are OpenAI/vendor claims even when they reference third-party benchmarks.
\end{itemize}
{\bfseries 読む際の境界}
\begin{itemize}[leftmargin=1.5em,itemsep=0.35em]
\item \textbf{分析上の留意点}: Benchmark, efficiency, and estimated-cost comparisons are largely vendor-reported and depend on effort levels, harnesses, token usage, and pricing assumptions; they should not be flattened into a universal ranking.
\end{itemize}
{\bfseries Primary source}
\begin{itemize}[leftmargin=1.5em,itemsep=0.25em]
\item \url{https://openai.com/index/advancing-the-price-performance-frontier-with-gpt-5-6}
\item \url{https://openai.com/index/gpt-5-6}
\item \url{https://openai.com/index/gpt-5-6-frontier-intelligence-efficiency}
\end{itemize}
{\scriptsize\color{SurveyMuted}Source-bound record: \texttt{evidence:SP-2026-M07:series-gpt-5-6-july-85e8f41ad0}.}
\end{technicalnote}
\medskip

\begin{technicalnote}{Claude Opus 5}{PRIMARY}
\begin{tabularx}{\linewidth}{@{}>{\bfseries}p{0.22\linewidth}X@{}}
Organization & Anthropic \\
Artifact type & MODEL \\
Chronology & 2026-07-24 (MODEL\_RELEASE) \\
\end{tabularx}
\smallskip
{\bfseries 一次資料から整理したtechnical points}
\begin{itemize}[leftmargin=1.5em,itemsep=0.35em]
\item \textbf{一次情報で確認できる事実}: Anthropic released Claude Opus 5 on July 24, 2026 and made it available across its platforms and API that day.
\item \textbf{一次情報で確認できる事実}: Anthropic priced Opus 5 at \$5 per million input tokens and \$25 per million output tokens and positioned it as the new default on Claude Max and the strongest model on Claude Pro.
\item \textbf{Vendor claim}: Anthropic reports Opus 5 as close to Fable 5 frontier capability at lower cost and presents broad coding, knowledge-work, science, and safety evaluation gains.
\end{itemize}
{\bfseries 読む際の境界}
\begin{itemize}[leftmargin=1.5em,itemsep=0.35em]
\item \textbf{分析上の留意点}: Benchmark and alignment comparisons are primarily Anthropic-reported; individual evaluation harnesses and partner testimonials are not equivalent to independent replication.
\end{itemize}
{\bfseries Primary source}
\begin{itemize}[leftmargin=1.5em,itemsep=0.25em]
\item \url{https://www.anthropic.com/news}
\end{itemize}
{\scriptsize\color{SurveyMuted}Source-bound record: \texttt{evidence:SP-2026-M07:item-official-index-anthropic-news-afddcb8068}.}
\end{technicalnote}
\medskip

\begin{technicalnote}{DeepSeek V4-Flash July API lifecycle}{PRIMARY}
\begin{tabularx}{\linewidth}{@{}>{\bfseries}p{0.22\linewidth}X@{}}
Organization & DeepSeek \\
Artifact type & MODEL\_UPDATE \\
Chronology & 2026-07-24 (API\_DEPRECATION); 2026-07-31 (MODEL\_UPDATE) \\
\end{tabularx}
\smallskip
{\bfseries 一次資料から整理したtechnical points}
\begin{itemize}[leftmargin=1.5em,itemsep=0.35em]
\item \textbf{一次情報で確認できる事実}: DeepSeek's changelog records the DeepSeek-V4-Flash API as entering public beta on July 31, 2026 under the model name `deepseek-v4-flash`.
\item \textbf{Vendor claim}: DeepSeek states that the July 31 build keeps the architecture and size of V4-Flash-Preview and changes post-training only.
\item \textbf{一次情報で確認できる事実}: An April 24 migration notice scheduled the legacy `deepseek-chat` and `deepseek-reasoner` names for discontinuation on July 24, 2026.
\end{itemize}
{\bfseries 読む際の境界}
\begin{itemize}[leftmargin=1.5em,itemsep=0.35em]
\item \textbf{分析上の留意点}: DeepSeek's benchmark scores for V4-Flash are vendor-reported and include a specified DeepSeek Harness configuration; this card treats the July lifecycle and API change as the primary fact.
\end{itemize}
{\bfseries Primary source}
\begin{itemize}[leftmargin=1.5em,itemsep=0.25em]
\item \url{https://api-docs.deepseek.com/updates/}
\end{itemize}
{\scriptsize\color{SurveyMuted}Source-bound record: \texttt{evidence:SP-2026-M07:item-official-index-deepseek-api-updates-72fd8b1ca6}.}
\end{technicalnote}
\medskip

\begin{technicalnote}{Gemini API July model lifecycle}{PRIMARY}
\begin{tabularx}{\linewidth}{@{}>{\bfseries}p{0.22\linewidth}X@{}}
Organization & Google \\
Artifact type & MODEL\_UPDATE \\
Chronology & 2026-07-21 (MODEL\_GA); 2026-07-30 (API\_PREVIEW) \\
\end{tabularx}
\smallskip
{\bfseries 一次資料から整理したtechnical points}
\begin{itemize}[leftmargin=1.5em,itemsep=0.35em]
\item \textbf{一次情報で確認できる事実}: Google made Gemini 3.6 Flash and Gemini 3.5 Flash-Lite generally available on July 21, 2026.
\item \textbf{一次情報で確認できる事実}: On July 30 Google added Gemini Robotics ER 2 and a streaming variant in public preview, with multimodal inputs and function calling for robotics workflows.
\item \textbf{Vendor claim}: Google describes Gemini 3.6 Flash as improving token efficiency and code/agentic planning while lowering price relative to 3.5 Flash.
\end{itemize}
{\bfseries 読む際の境界}
\begin{itemize}[leftmargin=1.5em,itemsep=0.35em]
\item \textbf{分析上の留意点}: Capability/efficiency positioning is Google-reported; the July release-notes page establishes model lifecycle and availability rather than independent benchmark superiority.
\end{itemize}
{\bfseries Primary source}
\begin{itemize}[leftmargin=1.5em,itemsep=0.25em]
\item \url{https://ai.google.dev/gemini-api/docs/changelog}
\end{itemize}
{\scriptsize\color{SurveyMuted}Source-bound record: \texttt{evidence:SP-2026-M07:item-official-index-google-gemini-api-re-d5b8ea25b2}.}
\end{technicalnote}
\medskip

\begin{technicalnote}{Kimi K3}{PRIMARY}
\begin{tabularx}{\linewidth}{@{}>{\bfseries}p{0.22\linewidth}X@{}}
Organization & Moonshot AI \\
Artifact type & OPEN\_WEIGHT \\
Chronology & 2026-07-16 (MODEL\_RELEASE) \\
\end{tabularx}
\smallskip
{\bfseries 一次資料から整理したtechnical points}
\begin{itemize}[leftmargin=1.5em,itemsep=0.35em]
\item \textbf{一次情報で確認できる事実}: Kimi Code's official changelog records Kimi K3 as released and open-sourced on July 16, 2026 and available in Kimi Code.
\item \textbf{Vendor claim}: Moonshot describes Kimi K3 as a 2.8-trillion-parameter model using KDA hybrid linear attention and Attention Residuals, with native visual understanding and up to a 1M-token context window.
\item \textbf{一次情報で確認できる事実}: The same July 16 Kimi Code release aligned the coder sub-agent with background tasks, plan mode, skills, and nested-agent capabilities.
\end{itemize}
{\bfseries 読む際の境界}
\begin{itemize}[leftmargin=1.5em,itemsep=0.35em]
\item \textbf{分析上の留意点}: Architecture scale, capability comparisons, and benchmark positioning are vendor claims from the Kimi documentation; open-source status should be interpreted using the actual published license/weights when drafting.
\end{itemize}
{\bfseries Primary source}
\begin{itemize}[leftmargin=1.5em,itemsep=0.25em]
\item \url{https://www.kimi.com/code/docs/en/kimi-code/whats-new.html}
\end{itemize}
{\scriptsize\color{SurveyMuted}Source-bound record: \texttt{evidence:SP-2026-M07:item-official-index-kimi-code-whats-new-788743b9bb}.}
\end{technicalnote}
\medskip
